\chapter*{Résumé}
\addcontentsline{toc}{chapter}{Résumé}
% Un \chapter* ne met pas le titre courant a jour : sans cette ligne, la seconde page du
% resume portait << Remerciements >> en tete.
\markboth{Résumé}{}

% LONGUEUR CONTRAINTE PAR LE § 5.1.c DES RECOMMANDATIONS JURY : 150 A 250 MOTS, pour le
% resume francais comme pour l'abstract anglais, ces deux textes servant a l'indexation sur
% le site de l'Institut. Mesure du 23 septembre 2026 AVANT resserrage : 322 et 305 mots sur
% le PDF, soit +29 % et +22 %. Ce qui a ete retire est du DETAIL, jamais un resultat ni une
% reserve : le compte des sources qui echouent et le champ VERIS rempli une fois sur 10 591,
% et le mecanisme de la decroissance de la charge avec la taille. Tous trois vivent dans la
% note de synthese et dans l'executive summary, qui n'ont AUCUNE limite de longueur.
% NE PAS REGROSSIR CES DEUX BLOCS : recompter sur le PDF apres toute retouche.
Ce mémoire quantifie le coût en capital d'un défaut de conformité au règlement DORA
(UE 2022/2554) pour un assureur exposé au risque cyber. La dépendance entre les cinq piliers y
est portée non par une copule mais par une cascade dirigée, normalisée à la Leontief de sorte que
la stabilité est \emph{garantie}, non supposée ; l'état de conformité pilote le capital par une
couche multi-états markovienne.

La contribution centrale n'est pas de \emph{mesurer} cette contagion mais d'en établir la
\emph{frontière d'identifiabilité} : sur données publiques agrégées, la \emph{direction} de la
propagation n'est pas identifiable quand sa \emph{co-occurrence} l'est. Plutôt que de poser une
matrice, le mémoire borne le capital sur l'ensemble des matrices admissibles : entre $6\,858$ et
$8\,697$~M\euro{} sur $1\,024$ sommets. Cette largeur de $1\,839$~M\euro{} est le coût de
l'ignorance, connu d'avance plutôt que subi.

Sur les quatre canaux que la conformité déplace, le besoin de capital sectoriel passe de
$6\,049$ à $20\,188$~M\euro, soit un facteur $3{,}34$. C'est un besoin au sens de l'ORSA, aucun
module DORA n'existant en Formule Standard. À l'échelle d'une entité notionnelle, entre $131{,}5$
et $183{,}3$~M\euro{} : un ordre de grandeur, le niveau absolu variant d'un facteur $41$ selon la
source. Résistent le rapport au socle et le classement des piliers.

La limite devient un livrable : un registre mieux spécifié resserrerait la bande de $66\,\%$, et
documenter une seule dépendance en lèverait $28{,}9\,\%$.

\paragraph{Mots-clés.} Risque cyber ; DORA ; SCR ; Solvabilité II ; théorie des valeurs
extrêmes ; cascade dirigée ; identification partielle ; bornes de capital ; inférence
bayésienne ; valeur de l'information.

\bigskip
\noindent\textbf{Abstract.} \emph{This thesis quantifies the capital cost of DORA non-compliance
(EU 2022/2554) for a cyber-exposed insurer. Dependence across the five DORA pillars is carried
not by a copula but by a directed cascade, Leontief-normalised so that stability is
\emph{guaranteed} rather than assumed. The firm's compliance state drives capital through a
Markov multi-state layer.}

\emph{The central contribution is not to \emph{measure} this contagion but to establish its
\emph{identifiability boundary}: from aggregated public data the \emph{direction} of propagation
is not identifiable, whereas its \emph{co-occurrence} is. Rather than positing a matrix, the
thesis bounds capital over the admissible set: between $6\,858$ and $8\,697$~M\euro{} across
$1\,024$ vertices. That width of $1\,839$~M\euro{} is the cost of not knowing, known in advance
rather than incurred.}

\emph{Across the four channels that compliance shifts, sector-level capital need rises from
$6\,049$ to $20\,188$~M\euro, a factor of $3.34$. This is an ORSA-type need, since the Standard
Formula has no DORA module. At the scale of a notional entity, between $131.5$ and
$183.3$~M\euro{}: an order of magnitude, the absolute level moving by a factor of $41$ with the
data source. What holds is the ratio to the baseline and the pillar ranking.}

\emph{The limitation becomes a deliverable: a better-specified incident register would narrow the
band by $66\,\%$, and documenting a single dependency would remove $28.9\,\%$ of the ambiguity.}

% =====================================================================
% USAGE DE L'IA GENERATIVE. EXIGE par le § 5.1.c des Recommandations Jury de
% l'Institut des actuaires (1er avril 2026) : << En complement et a la fin du
% resume et distinct de ce dernier, quelques phrases concises mentionneront
% l'usage effectue de l'intelligence artificielle generative tant au cours des
% travaux que dans la redaction du memoire. >> Le § 4.5 precise que tout
% manquement a la transparence conduit a l'ajournement sans autre motivation.
%
% KELIAN : A RELIRE ET A AJUSTER AVANT LE DEPOT. Ce paragraphe decrit l'usage
% tel que le depot en porte la trace. Si le perimetre reel est plus large ou
% plus etroit, c'est ce paragraphe et l'annexe d'inventaire qui doivent bouger, pas le
% principe : le referentiel demande une transparence TOTALE, et une declaration
% incomplete est traitee comme une absence de declaration.
%
% AUCUN CHIFFRE ICI A DESSEIN : un nombre pose dans ce bloc irait chercher sa
% confirmation dans les scripts cites en bas de fichier, qui ne le portent pas.
% Les comptes du dispositif de verification vivent dans l'annexe, qui se declare
% hors script pour ce motif.
% =====================================================================
\bigskip
\noindent\textbf{Usage de l'intelligence artificielle générative.} Ce mémoire a été rédigé avec
l'appui ponctuel d'outils d'IA générative, utilisés comme assistance à la rédaction et à la reformulation du
manuscrit, à l'écriture des scripts d'analyse et à la vérification du texte, sous le contrôle
permanent de l'auteur. L'ensemble des analyses, choix méthodologiques et interprétations des
résultats relèvent exclusivement de l'auteur, et aucun résultat chiffré n'est produit par
l'assistant : tout nombre publié provient de l'exécution d'un script du dépôt. L'annexe~\ref{ann:ia}
détaille cet usage ainsi que les mesures prises pour en contrôler la validité.

% =====================================================================
% SOMMAIRE AU NIVEAU DES SECTIONS, et non des sous-sections. Mesure du
% 22 septembre 2026 : en listant les sous-sections il occupait SIX pages, sur un
% document que Kelian veut imperativement sous 200. Rien n'est retire du
% memoire, seule la profondeur du sommaire change ; les sous-sections gardent
% leur numero et leur titre dans le corps.
\setcounter{tocdepth}{1}
\tableofcontents

% =====================================================================

% SOURCES-SCRIPTS: 28 39 33 43 68 92 66 89 91 100 30 55 65
